\section*{Agents4Science AI Involvement Checklist}

This checklist documents the role of AI in our research across different aspects of the scientific process. We provide scores and explanations for each category to ensure transparency about AI involvement.

\begin{enumerate}
    \item \textbf{Hypothesis development}: Hypothesis development includes the process by which you came to explore this research topic and research question. This can involve the background research performed by either researchers or by AI. This can also involve whether the idea was proposed by researchers or by AI. 

    Answer: \textbf{Mostly AI, assisted by human}
    
    Explanation: The core hypothesis that AI-driven ensemble approaches could effectively optimize OpenMP parallelization for GraphViz algorithms was primarily developed through AI analysis of existing literature and identification of research gaps. AI systems analyzed performance bottlenecks and proposed the multi-model ensemble strategy. Human researchers provided domain expertise and guided the focus toward GraphViz and Apple M1 architecture.
    
    \item \textbf{Experimental design and implementation}: This category includes design of experiments that are used to test the hypotheses, coding and implementation of computational methods, and the execution of these experiments. 

    Answer: \textbf{Mostly AI, assisted by human}
    
    Explanation: AI systems designed the comprehensive experimental framework, including the multi-level performance measurement methodology, statistical validation protocols, and the AI ensemble architecture. AI generated the OpenMP optimization code and implemented the profiling pipeline. Human researchers provided oversight for experimental validity, hardware configuration, and ensured adherence to scientific standards.
    
    \item \textbf{Analysis of data and interpretation of results}: This category encompasses any process to organize and process data for the experiments in the paper. It also includes interpretations of the results of the study.
 
    Answer: \textbf{Mostly AI, assisted by human}
    
    Explanation: AI systems performed the majority of data analysis, including statistical computations, performance trend identification, and interpretation of optimization effectiveness. AI generated the comprehensive performance visualizations and identified key insights about thread efficiency and scalability patterns. Human researchers validated the statistical methodology and provided domain-specific interpretation of results.
    
    \item \textbf{Writing}: This includes any processes for compiling results, methods, etc. into the final paper form. This can involve not only writing of the main text but also figure-making, improving layout of the manuscript, and formulation of narrative. 

    Answer: \textbf{AI-generated}
    
    Explanation: The paper was primarily written by AI systems, including the technical content, methodology descriptions, results analysis, and narrative structure. AI generated all figures, tables, and visualizations. AI also handled the literature review, citation management, and formatting. Human researchers provided minimal guidance on structure and ensured compliance with conference requirements.

    \item \textbf{Observed AI Limitations}: What limitations have you found when using AI as a partner or lead author? 

    Description: Key limitations observed include: (1) Occasional inconsistencies in technical details that required human verification, (2) Tendency to over-optimize prose that sometimes obscured clarity, (3) Challenges in maintaining consistent notation across complex technical sections, (4) Difficulty in balancing comprehensive coverage with conciseness constraints, and (5) Need for human oversight to ensure experimental protocols met rigorous scientific standards. Despite these limitations, the AI ensemble approach significantly accelerated research productivity while maintaining high technical quality.
\end{enumerate}
